% !TEX root = main_memo.tex

\chapter{Some known and less known preliminaries}\label{sectionPrelim}

Given a function $f:X\to Y$, we use the following classical notations. We denote by $\dom f=X$ the domain of $f$ and refer to $Y$ as the codomain of $f$.
We write $\im f =\set{f(x)}[x\in X]$ for the image, or the range, of $f$. 
For $A\subseteq X$, $f\restr{A}:A\to Y$\index[IdxSymb]{f@{$f\restr{B}$}} stands for the restriction of $f$ to $A$, and $\id_{X}:X\to X$\index[IdxSymb]{id@{$\id_X$}} denotes the identity function on $X$.
For $B\subseteq Y$, $f^{-1}(B)=\set{x\in X}[f(x)\in B]$ denotes the preimage of $B$ by $f$, and we call \emph{co-restriction}\index[IdxNotion]{co-restriction} of $f$ to $B$ the function $f\restr{f^{-1}(B)}$, denoted by $f\corestr{B}$\index[IdxSymb]{f@{$f\corestr{B}$}}.
%Given a function $f$ and a set $X$, we use the following classical notations:
%$\im(f)$ stands for the image, or the range, of $f$, $f^{-1}(X)$ denotes the preimage of $X$ by $f$,
%$f\restr{X}$ stands for the restriction of $f$ to $X$, and $\id_{X}$ denotes the identity function on $X$.
To improve readability, we often use the multiplicative notation $fg$ for the composition $f\circ g$ of $g$ with $f$.


When a pair $(\sigma,\tau)$ of continuous functions witnesses that a function $f$ reduces continuously to a functions $g$, we write $f\leq g$ and we say that $(\sigma,\tau)$ \emph{continuously reduces}, or simply \emph{reduces}, $f$ to $g$.
As usual with quasi-orders, there is an induced equivalence relation: we write $f\equiv g$ when both $f\leq g$ and $g\leq f$ hold,
and we say that $f$ and $g$ are \emph{continuously equivalent}.
We also denote by $f<g$ the relation $f\leq g$ and $f\not\geq g$, and we say that $f$ and $g$ are \emph{incomparable} when both $f\not\leq g$ and $g\not\leq f$. \index[IdxSymb]{0@{$f\leq g$, $f<g$, $f\equiv g$}}
An \emph{antichain} is a set of pairwise incomparable elements, and an \emph{infinite strictly descending chain} is a sequence $(f_n)_{n\in\N}$
satisfying $f_{n+1}<f_n$ for all $n\in\N$. 




%Given a function $f:X\to Y$ and $B\subseteq Y$, we call \emph{co-restriction} of $f$ to $B$ the function $f\restr{f^{-1}(B)}$ and we denote it by $f\corestr{B}$.

It follows from the definition of continuous reducibility that any function $f:X\to Y$ is continuously equivalent to the surjective function $f\corestr{\im f}:X\to \im f$,
where $\im f$ is endowed with the subspace topology\footnote{However, restricting our attention to surjective functions would not simplify proofs nor notations.
Indeed some functions like the centered functions from \cref{sectionCentered} are equivalent to their corestriction to arbitrary small subset of their image.}. 

\smallskip

About the existence of a reduction $(\sigma,\tau)$ between functions $f$ and $g$, it is useful to notice that to any given (potential) $\sigma$ can only correspond to a single $\tau$.
More precisely, $f\leq g$ if and only if there exists a continuous $\sigma:X\to X'$ such that the relation  
\[
T=\set{g\sigma(x), f(x))}[x\in X]
\]
is the graph of a continuous function $\tau: \im (g\sigma) \to \im(f)$.

While $\sigma$ fully determines $\tau$, a given candidate $\tau$ can potentially correspond to various $\sigma$.
Indeed, $f\leq g$ if and only if there exists a subset $B$ of $Y'$ and a continuous surjection $\tau:B\to \im(f)$ such that the relation 
\[
\set{(x,y)\in X\times X'}[g(y)\in B \text{ and } \tau g(y) =f(x)]
\]
admits a continuous uniformization (for more on uniformizations, see \cite[Section 18]{kechris}).

\medskip

We next show that continuous reducibility between functions respects measurability and generalizes both (coarse) Wadge reducibility and embeddability.
\index[IdxNotion]{continuous reduction!Wadge reducibility}\index[IdxNotion]{topological embeddability!between spaces}

Recall that for two topological spaces $X$ and $Y$, $A\subseteq X$ \emph{Wadge reduces to} $B\subseteq Y$ if there is a continuous function $\sigma:X\to Y$
such that $\sigma^{-1}(B)=A$, denoted by $A\leq^{X,Y}_WB$\index[IdxSymb]{2@{$\leq^{X,Y}_W$}}.
Given $A\subseteq X$ the notation $\1_{A,X}:X\to \{0,1\}$\index[IdxSymb]{one@{$\1_{A,X}$}} stands for the characteristic function of $A$ in $X$.

A continuous function $\sigma:X\to Y$ is a topological embedding (or simply an embedding) if and only if there is a continuous function $\tau:\im(\sigma)\to X$
such that $\tau\sigma=\id_X$, we also say that $(\sigma,\tau)$ \emph{witnesses} the embedding from $X$ to $Y$.

We sometimes write \enquote{reduces} and \enquote{embeds} instead of \enquote{continuously reduces} and \enquote{topologically embeds}.

Finally, a \emph{pointclass}\index[IdxNotion]{pointclass} is a functional class $\mathbf{\Gamma}$%\index[IdxSymb]{$\mathbf{\Gamma}$}
 mapping any topological space $X$ to a subset of the power set of $X$ with the property
that given any continuous function $\sigma:X\to Y$ if $A\in\mathbf{\Gamma}(Y)$ then $\sigma^{-1}(A)\in\mathbf{\Gamma}(X)$.
Observe that $\mathbf{\Gamma}(X)$ is closed under preimages by continuous functions from $X$ to $X$, for all spaces $X$.
A function $f:X\to Y$ is said \emph{$\mathbf{\Gamma}$-measurable} if the preimage of an open set of $Y$ by $f$ is in $\mathbf{\Gamma}(X)$.

\begin{proposition}\label{LinkswithWadgeEmbedMeas}
Let $X,X',Y$ and $Y'$ be topological spaces.
\begin{enumerate}
\item $X$ embeds in $Y$ if and only if $\id_X\leq\id_Y$.
\item If $A\subseteq X$ and $B\subseteq Y$ then $\1_{A,X}\leq\1_{B,Y}$ if and only if $A\leq^{X,Y}_WB$ or $A\leq^{X,Y}_WY\setminus B$.
\item Given a pointclass $\mathbf{\Gamma}$, if $f:X\to Y$ reduces to $g:X'\to Y'$ and $g$ is $\mathbf{\Gamma}$-measurable, then so is $f$.
\end{enumerate}
\end{proposition}


\begin{proof}
A space $X$ embeds in $Y$ if and only if a pair of maps $(\sigma,\tau)$ witnesses the embedding from $X$ to $Y$ if and only if $(\sigma,\tau)$ reduces $\id_X$ to $\id_Y$.

There are only two (continuous) functions $\set{0,1}\to\set{0,1}$, the identity and the function swapping $0$ and $1$.
So, given $A\subseteq X$ and $B\subseteq Y$, $(\sigma,\tau)$ reduces $\1_{A,X}$ to $\1_{B,Y}$ if and only if $A\leq^{X,Y}_WB$ (in case $\tau$ is the identity)
or $A\leq^{X,Y}_WY\setminus B$ (in case $\tau$ swaps $0$ and $1$).

Fix a continuous reduction $(\sigma,\tau)$ from $f$ to $g$, and assume that $U\subseteq Y$ is open. Note that we have $f^{-1}(U)=(g\sigma)^{-1}(\tau^{-1}(U))=\sigma^{-1}(g^{-1}(\tau^{-1}(U)))$. By continuity of $\tau$, $\tau^{-1}(U)$ is open in $\im (g\sigma)$ so there exists an open set $V$ in $Y'$ with $\tau^{-1}(U)=V\cap \im (g\sigma)$. Since $g$ is $\mathbf{\Gamma}$-measurable and $\sigma$ is continuous, it follows that $g\sigma:X\to Y'$ is $\mathbf{\Gamma}$-measurable. Therefore, we get that $f^{-1}(U)=(g\sigma)^{-1}(\tau^{-1}(U))=(g\sigma)^{-1}(V)\in \mathbf{\Gamma}(X)$, as desired.
\end{proof}

We use the following basic facts all along the paper without reference.

\begin{proposition}
Let $f$ be a function between topological spaces.
\begin{enumerate}
\item If $X\subseteq\dom f$ then $f\restr{X}\leq f$,
\item if $f$ is continuous and $\dom f$ embeds in $X$ then $f\leq\id_X$, and
\item if $f$ is continuous and $\im f$ embeds in $Y$ then $f\leq\id_Y$.
\end{enumerate}
\end{proposition}

\begin{proof}
The desired continuous reductions are:
$(\id_X, \id_{f(X)})$; $(\sigma, f\tau)$ with $(\sigma,\tau)$ witnessing the embedding from $\dom f$ into $X$;
and finally $(\sigma f, \tau)$ with $(\sigma,\tau)$ witnessing the embedding from $\im f$ into $Y$.
\end{proof}

Using the universal properties of $\Q$, $\cantor$ and $\cN$ \cite[(7.12),(7.8)]{kechris}, we have the following corollary.
\begin{Corollary}\label{ConsequencesofUniversality}
Let $f:X\to Y$ be continuous.
\begin{enumerate}
\item If $\im f$ is countable and metrizable, then $f\leq \id_\Q$.
\item If $X$ is zero-dimensional separable metrizable, then $f\leq \id_{\cN}$ and $f\leq \id_{\cantor}$. 
\end{enumerate}
In particular, $\id_{2^\N}\equiv \id_{\cN}$.
%\begin{enumerate}
%\item $\id_{2^\N}\equiv \id_{\cN}$.
%\item If $f$ is continuous with countable metrizable image, then $f\leq \id_\Q$.
%\item If $f$ is continuous with zero-dimensional separable metrizable domain, then $f\leq \id_{\cN}$. 
%\end{enumerate}
\end{Corollary}


Finally, recall from the introduction that when a reduction $(\sigma,\tau)$ from $f$ to $g$ consists of topological embeddings, we say that $f$ \emph{embeds topologically}\footnote{Observe that this actually coincides with   the following definition used in the articles cited in the introduction: $f$ embeds in $g$ if there is a pair $(\sigma, \tau)$ of embeddings such that $\tau f=g\sigma$. See \cite[Section 3]{embeddabilityCPZ} for more details.}
 in $g$, we denote it by $f\segm g$\index[IdxSymb]{0@{$f\segm g$}}.\index[IdxNotion]{topological embeddability!between functions}
Topological embeddability and continuous reducibility coincide on topological embeddings.

\begin{proposition}\label{ContRedonEmbed}
Let $f$ and $g$ be functions between topological spaces.
Assume that $(\sigma,\tau)$ reduces $f$ to $g$.

If $f$ is injective, then $\sigma$, $g\restr{\im \sigma}$ and $\tau$ are injective too.

If moreover $f$ is an embedding, then so is $g\restr{\im \sigma}$, and actually $(\sigma,\tau)$ embeds $f$ in $g$.
\end{proposition}

\begin{proof}
To see the first point, note that if any of those maps is not injective, nor is $f$ since $f=\tau g \sigma$ holds.

If now $f$ is an embedding, by definition there is an embedding $h$ such that $hf=\id_{\dom f}$, but since $hf= h\tau g\sigma=\id_{\dom f}$, we know that $\sigma$ is an embedding too.
Similarly $fh=\tau g\sigma h=\id_{\im f }$ and $\tau$ is an embedding. Finally we get $g\sigma h\tau=\id_{\im \sigma}$, which concludes.
\end{proof}



We say that a function $f:X\to Y$ is \emph{homeomorphic} to a function $f':X'\to Y'$ if $f\leq f'$ is witnessed by $(\sigma,\tau)$ where $\sigma:X\to X'$ and $\tau:\im f'\to \im f$ are both homeomorphisms.\index[IdxNotion]{homeomorphism between functions}
Note that in this special case, not only does $f$ embed in $f'$ and reduces to $f'$, but also the pair $(\sigma^{-1},\tau^{-1})$ is a homeomorphism from $f'$ to $f$. Let us denote by $\functionsonbaire$ the class of functions $f:A\to B$ where $A,B$ are subsets of the Baire space $\cN$.\index[IdxSymb]{F@{$\functionsonbaire$}}

\begin{proposition}\label{RepresentationforFunctions}
Any function between $0$-dimensional separable metrizable spaces is homeomorphic to a function in $\functionsonbaire$. 
So is any function from a $0$-dimensional separable metrizable space to a countable metrizable space. 
\end{proposition}

\begin{proof}
Because any $0$-dimen\-sional separable metrizable space is homeomorphic to a subset of $\cN$ and every countable metrizable space is homeomorphic to a subset of $\Q$, hence $0$-dimensional.
\end{proof}



\section{Non-scattered functions}

Recall that a topological space $X$ is \emph{scattered} if any of its non-empty subset has an isolated point.\index[IdxNotion]{space!scattered}\index[IdxNotion]{scattered!space}
For a metrizable space, $X$ is scattered if and only if $\Q$ does not embed into $X$ (see for instance \cite[Théorème 2]{MR0107849}, or \cite[Notes on (7.12)]{kechris}).
We show an analogous characterization for continuous functions.
\begin{definition}
We say that a function $f$ is \emph{scattered}\index[IdxNotion]{space!scattered}\index[IdxNotion]{scattered!space} if for every non-empty subset $A\subseteq \dom f$ the restriction $f\restr{A}$ is constant on some non-empty open subset of $A$.
\end{definition}
%Recall that we say that a function $f$ is \emph{scattered} if 
% for every non-empty subset $A\subseteq \dom f$ the restriction $f\restr{A}$ is constant on some non-empty open subset of $A$.
 Note that $f$ is not scattered if and only if it admits a non-empty restriction $f\restr{A}$ which is \emph{nowhere locally constant},
 that is such that $f\restr{U}$ is not constant for every non-empty open subset $U$ of $A$. 
Remark that  the identity on the space of rationals $\id_\Q$ is not scattered. For continuous functions between metrizable spaces, this function is actually the only obstruction to being scattered.
%We will see that for any function $f$, if $\id_\Q\leq f$ then $f$ is not scattered (\cref{CBbasicsfromJSL}).
%For continuous functions between metrizable spaces, the identity on the space of rationals is actually the only obstruction to being scattered:


\begin{theorem}\label{prop:nlc_implies_nonscattered}
Let $f:X\to Y$ be a continuous function from a metrizable space to a Hausdorff space.
If $f$ is not scattered, then $\id_\Q\segm f$.  
\end{theorem}

\begin{proof}
If $f$ is not scattered, then it admits a non-empty restriction $f\restr{A}$ that is nowhere locally constant and $f\restr{A}\segm f$.
Hence, it suffices to show that if $f$ is nowhere locally constant, then $\id_\Q\segm  f$.

Let $d$ be a compatible metric for $X$.
We make the following observation: if $x\in X$, $\epsilon>0$ and $f(B(x,\epsilon))\subseteq U$ for some open set $U$ of $Y$,
then there exists $x'\in B(x,\epsilon)$, $\epsilon'<\epsilon$ and open sets $U_0,U_1\subseteq U$ such that $B(x,\epsilon'),B(x',\epsilon')\subseteq B(x,\epsilon)$,
$U_0\cap U_1=\emptyset$, $f(B(x,\epsilon'))\subseteq U_0$ and $f(B(x',\epsilon'))\subseteq U_1$.
To see this, note that since $f$ is nowhere locally constant there exists $x'\in B(x,\epsilon)$ such that $f(x)\neq f(x')$.
Since $Y$ is Hausdorff, there exists disjoint neighborhoods $U_0,U_1\subseteq U$ separating $f(x)$ from $f(x')$.
By continuity of $f$, we can find a small enough $\epsilon'<\epsilon$ such that $f(B(x,\epsilon'))\subseteq U_0$ and $f(B(x',\epsilon'))\subseteq U_1$ and $B(x,\epsilon'), B(x',\epsilon')\subseteq B(x,\epsilon)$.

We use this to define Cantor schemes\footnote{See \cite[Definition 6.1]{kechris} for a definition of \emph{Cantor scheme}.}
$(C_s)_{s\in 2^{<\N}}$ consisting of open balls $B(x_s, \rho_s)$ of $X$ and $(U_s)_{s\in 2^{<\N}}$ consisting of open sets of $Y$ such that for all $s\in 2^{<\N}$:
\begin{enumerate}
\item $f(C_{s\conc(i)})\subseteq U_{s\conc(i)}$ for $i=0,1$,
\item $\rho_s\leq 2^{-|s|}$.
\end{enumerate}
Choose $x_\emptyset\in X$, let $\rho_\emptyset=1$, set $C_\emptyset=B(x,1)$ and $U_\emptyset=Y$. Assume we have defined $C_s=B(x_s, \rho_s)$.
The preliminary observation allows us to find $x'$ and $\rho\leq 2^{-(|s|+1)}$ such that $B(x_s,\rho)$ and $B(x',\rho)$ are included in $C_s$ and whose images by $f$ are separated by open sets $U_0,U_1\subseteq U_s$.
We set $x_{s\conc(0)}=x_s$, $x_{s\conc(1)}=x'$, $\rho_{s\conc(0)}=\rho_{s\conc(1)}=\rho$ and $U_{s\conc (i)}=U_i$ for $i=0,1$. 

Note that by construction, we have $\{x_{s}\}= \bigcap_{n\in\N}C_{s\conc (0)^n}$ for every $s\in 2^{<\N}$.
Identifying $\Q$ with the subset of $\cantor$ of sequences that are eventually constant equal to $0$, our Cantor scheme induces a function $\sigma:\Q\to X$ which is an embedding by \cite[(7.6)]{kechris}. 
Moreover, $f\sigma$ is an embedding too. 
To see this, note that $f\sigma(\Q\cap N_s)=f\sigma(\Q)\cap U_s$ for all $s\in 2^{<\N}$.
Indeed if $x\in \Q\cap N_s$ then by construction $\sigma(x)\in C_s$ and so $f\sigma(x)\in U_s$.
Conversely, if $x\notin \Q\cap N_s$ then $x\in N_t$ for some $t\in 2^{<\N}$ incompatible with $s$ and so $f\sigma(x)\in U_t$ but then $f\sigma(x)\notin U_s$ since $U_s\cap U_t=\emptyset$.

Since $f\sigma$ is an embedding, its inverse $\tau:\im (f\sigma)\to \Q$, $f\sigma(x)\mapsto x$ is an embedding too and clearly $\id_\Q=\tau f \sigma$. Therefore $\id_\Q\segm f$, as desired.
\end{proof}

%Combining \cref{prop:nlc_implies_nonscattered} with \cref{ConsequencesofUniversality} gives
%
%\begin{corollary}\label{nonscatteredfunctions}
%Any non-scattered continuous function from a metrizable space to a countable metrizable space is continuously equivalent to $\id_\Q$. 
%\end{corollary}

Continuity\footnote{Metrizability as well, see \cite[Footnote of Théorème 2]{MR0107849}.} is needed in \cref{prop:nlc_implies_nonscattered}
\begin{example}[A non-scattered function that does not reduce $\id_\Q$]\label{counterexamplenonscat}
Let $f$ be the characteristic function of a dense and co-dense subset of $\Q$.
Then $f$ is not scattered because it is nowhere locally constant, as its image on every basic open set is exactly $2$.
Being discontinuous and having finite image, it cannot continuously reduce $\id_\Q$.
\end{example}

For a finite basis result for all functions from $\Q$ to an analytic metric space with respect to closed topological embeddability, see \cite[Theorem 5.5]{carroymiller2}.

\medskip

Our next goal is to understand the continuous functions with uncountable range up to continuous reducibility.
For a continuous function $f$ with a Polish domain and separable metric image, the proof of \cref{prop:nlc_implies_nonscattered} can be adapted to show that if $f$ is not scattered,
then the constructed Cantor schemes are actually witnessing $\id_{\cantor}\segm f$.

We next show that we can extend this phenomenon to analytic domains and Borel functions,
for which we first establish equivalent conditions for the embeddability of $\id_{\cantor}$ into a function.
Recall that a function is \emph{Baire-measurable} if preimages of open sets have the Baire property\index[IdxNotion]{Baire-measurable},
and that $\Pee{2}$ sets are countable intersections of open sets.

\begin{proposition}\label{EquivEmbedCantor}
Let $f$ be a function between topological spaces. The following are equivalent:
\begin{enumerate}
\item There is an embedding $\sigma:\cantor\to\dom(f)$ such that $f\sigma$ is an embedding.
\item There is a Baire-measurable $\sigma:\cantor\to\dom(f)$ such that $f\sigma$ is an embedding.
\item There is $K\subseteq \dom(f)$ homeomorphic to $\cantor$ such that $f\restr{K}$ is both injective and Baire-measurable.
\end{enumerate}
\end{proposition}

\begin{proof}
Any embedding -- being continuous -- is Baire-measurable, so the first item implies the second one.

%Any $\sigma$ such that $f\sigma$ is injective must be injective, and if moreover such $\sigma$ is Baire-measurable on $\cantor$ 
If $f\sigma$ is injective, then so is $\sigma$. Since $\sigma$ is moreover Baire-measurable on $\cantor$ 
then by \cite[(8.38)]{kechris} it must be continuous on a dense $\Pee{2}$ subset $G$ of $\cantor$.
In particular, $G$ is a perfect Polish space and so there exists an embedding $\sigma':\cantor \to G$ \cite[(6.2)]{kechris}.
%Now $\sigma\sigma'$ is both injective and continuous, hence an embedding by compactness of $\cantor$, so $f\sigma\sigma'$ is also an embedding and we obtained the first item.
Now $\sigma\sigma'$ is both injective and continuous, so $K=\im \sigma\sigma'$ is homeomorphic to $\cantor$ by compactness. Moreover, $f\restr{K}=(f\sigma) \sigma^{-1}$ is an embedding as a composition of embeddings. Hence $f\restr{K}$ is injective and Baire measurable, which shows that the second item implies the third one.

%Now if $\sigma$ is an embedding and $f\sigma$ is an embedding, then $f\restr{\im(\sigma)}$ is also an embedding, in particular it is injective and Baire-measurable, so we have the third item.
Finally, suppose there is $K\subseteq \dom(f)$ homeomorphic to $\cantor$ such that $f\restr{K}$ is both injective and Baire-measurable. By
\cite[(8.38)]{kechris} again, $f\restr{K}$ is continuous on a dense $\Pee{2}$ subset of $K$. So there exists an embedding $\sigma:\cantor\to K$
such that $f\sigma$ is both injective and continuous, hence an embedding by compactness% of $\cantor$
, and so the first item holds.
\end{proof}

We say that a class of functions has the \emph{Perfect Function Property}\index[IdxNotion]{Perfect Function Property} if every function in it either has countable image or $\id_{\cantor}$ embeds in $f$.
Observe that a class of spaces has the perfect set property if and only if the class of their identity functions has the perfect function property.

The following proposition has various proofs. It was first proven in \cite[Proposition 2.1]{carroy2013quasi} using Silver's Theorem and
 a game-theoretic proof was recently found by Lutz and Siskind \cite{lutz20231}.
 These proofs both use the third item from \cref{EquivEmbedCantor}; here we propose another proof, suggested to us by Zolt\'an Vidny\'anszky.

A subset $A$ of a topological space $X$ is \emph{Borel} if it is in the $\sigma$-algebra generated by the open sets,
and a function is \emph{Borel} if it is Borel-measurable. Note that zero-dimensionality is not used in the following result.

\begin{proposition}\label{uncountablerange}
The class of Borel functions from analytic to separable metrizable spaces has the Perfect Function Property.
\end{proposition}

\begin{proof}
Let $f:X\to Y$ be Borel with uncountable image. As $X$ is analytic we can take a continuous surjection $\phi:\cN\to X$.
Since $\im(f\phi)$ is analytic and its completion is an uncountable Polish space, the Perfect Set Theorem for analytic sets \cite[(29.1)]{kechris} ensures that there is an embedding $\tau:\cantor\to\im(f\phi)$.
The relation $xRx'$ if and only if $\tau(x)=f\phi(x')$ is analytic in $\cantor\times\cN$, so by the Jankov--von Neumann uniformization Theorem (\cite[(29.9)]{kechris})
it has a Baire-measurable uniformizing function $\sigma:\cantor\to\cN$, so $f\phi\sigma=\tau$ is an embedding, and this concludes by \cref{EquivEmbedCantor}.
\end{proof}


\section{Scattered functions}\label{sectionCB}

We generalize the Cantor--Bendixson analysis of spaces to functions, which was already initiated in~\cite{carroy2013quasi}. % the basics of the Cantor--Bendixson analysis for functions.

\index[IdxNotion]{Cantor--Bendixson!analysis}
Given a function $f$ and $x\in\dom f$, we say that $x$ is \textit{$f$-isolated} if $f^{-1}(\{f(x)\})$ is a neighbourhood of $x$, \ie $f$ is locally constant at $x$.
Note that a function $f$ is scattered if and only if every non-empty subset $A$ of $\dom f$ contains an $f\restr{A}$-isolated point.
Given $A\subseteq\dom f$, we let $I(f,A)$ stand for the set of all $f\restr{A}$-isolated points. It is an open subset of $A$.
We now define inductively the following closed sets of $\dom f$:\index[IdxNotion]{Cantor--Bendixson!derivative}
\begin{align*}
\CB_0(f)&=\dom f,\\
\CB_{\alpha+1}(f)&=\CB_\alpha(f)\backslash I(f,\CB_\alpha(f)),\\ 
\CB_\lambda(f)&=\bigcap_{\alpha<\lambda}\CB_\alpha(f)\mbox{ for }\lambda\mbox{ limit.}
\end{align*}
\index[IdxSymb]{C@{$\CB_\alpha(f)$}}\index[IdxSymb]{C@{$\CB(f)$}}
We set $\CB(f)=\min\set{\alpha}[\CB_\alpha(f)=\CB_{\alpha+1}(f)]$, and
we call it the \textit{Cantor--Bendixson rank} of $f$, or \emph{$\CB$-rank} for short.\index[IdxNotion]{Cantor--Bendixson!rank}
Recall from the introduction that $\ker_{\CB}{f}:=\CB_{\CB(f)}(f)$ is called the perfect kernel of $f$, and $f$ is nowhere locally constant on it.
In particular, functions with an empty perfect kernel play a central role in this paper as are exactly the scattered functions. 
Note that this equivalence stated in introduction as part of \cref{scatterediffemptykernel} actually holds for functions between topological spaces. \index[IdxNotion]{scattered!function}\index[IdxNotion]{function!scattered}

\begin{proposition}\label{scatterediffemptykernel_general}
Let $f$ and $g$ be two functions between topological spaces. Then $f$ is scattered if and only if $f$ has empty perfect kernel. 
\end{proposition}

\begin{proof}
If $f$ has non-empty perfect kernel $A$, then by definition of the Cantor--Bendixson derivative $f\restr{A}$ is nowhere locally constant, so $f$ is not scattered.
Conversely, if $f$ has empty perfect kernel, then for some ordinal $\alpha$ we have $\dom(f)=\bigcup_{\beta<\alpha}I(f,\CB_\beta(f))$,
so given any non-empty $A\subseteq\dom(f)$, setting $\gamma=\min\set{\beta<\alpha}[A\cap I(f,\CB_\beta(f))\neq\emptyset]$,
then $A\subseteq \CB_\gamma(f)$ and $f\restr{A}$ is locally constant at any $x\in A\cap I(f,\CB_\gamma(f))$, which shows that $f$ is scattered.
\end{proof}
The Cantor--Bendixson rank is defined using a closed derivative, so any function with a second countable domain has a countable $\CB$-rank\footnote{See \cite[Theorem 6.9]{kechris}}.
Since a locally constant function on a second countable space has countable range, it follows from \cref{scatterediffemptykernel_general} that any scattered function with second countable domain has countable image. 
In particular, every scattered continuous functions from a zero-dimensional separable metrizable space to a metrizable space, has countable range and therefore is homeomorphic to a function in $\functionsonbaire$ by \cref{RepresentationforFunctions}. With this in mind, we define the class of functions which is the main focus of this article as follows: 
{
%\begin{definition}
%The class $\sC$ denotes the class of functions from a separable metrizable zero-dimen\-sional space to a metrizable space which are both continuous and scattered. 
%\end{definition}
\begin{definition}
The class $\sC$\index[IdxSymb]{Scat@{$\sC$}} denotes the class of functions $f\in \functionsonbaire$ which are continuous and scattered. 
\end{definition}
 }
We are now in a position to show that \cref{Introthm:BQO} reduces to \cref{Introthm:BQOonScat}.
Recall that both \wqo{} and \bqo{} are closed under finite unions, so if $Q$ and $Q'$ are \wqo{} (resp. \bqo{}) then so is $Q\cup Q'$.\index[IdxNotion]{better-quasi-order (\bqo{})}
Recall as well that quasi-orders whose induced equivalence relation has finitely many classes are \bqo{} (see \cite[Example 3.8 and Proposition 3.11]{wellbetterinbetween}).
As announced we have:

\begin{theorem}\label{FirststepforBQOthm} 
For every function $f$ that satisifies the hypothesis of \cref{Introthm:BQO}, one of the following holds: $f\equiv \id_\Q$, $f\equiv \id_{\cN}$ or $f$ is scattered.
In particular, \cref{Introthm:BQOonScat} implies \cref{Introthm:BQO}.
\end{theorem}

\begin{proof}
Let $f:X\to Y$ be a continuous function with $X$ zero-dimensional separable metrizable and $Y$ metrizable. If $f$ is scattered then being continuous with second countable domain it has countable range. So in light of \cref{RepresentationforFunctions}, there exists $\hat{f}\in\sC$ such that $f$ is homeomorphic to $\hat{f}$ and in particular $f\equiv \hat{f}$.

So suppose that $f$ is non scattered, so $\id_\Q\segm f$ by \cref{prop:nlc_implies_nonscattered}. If $\im f$ is furthemore countable, then actually $f\equiv \id_\Q$ by \cref{ConsequencesofUniversality}. Otherwise if $X$ is furthermore analytic, then $f$ has the Perfect Function Property by \cref{uncountablerange} and so $\id_{\cN}\segm f$. In this case, since $X$ is zero-dimen\-sional separable metrizable, then in fact $f\equiv \id_{\cN}$ \cref{ConsequencesofUniversality}.
% Using \cref{nonscatteredfunctions,nonscatteredfunctions_PFP,uncountablerange}, we see that any non-scattered function that satisfies the hypotheses of \cref{Introthm:BQO} is equivalent to either $\id_\Q$ or $\id_{\cN}$.

Therefore, up to continuous equivalence, the class of functions described in \cref{Introthm:BQO} can be written as the union of $\sC$ together with the two equivalence classes of $\id_\Q$ or $\id_{\cN}$. Hence if $\sC$ is \bqo{} then the whole class is \bqo{}.
\end{proof}

%
%
%\begin{aynote}
%Dans un sens, je pense que ce serait même mieux de définir $\sC\subseteq\functionsonbaire$. Cela est possible par \cref{RepresentationforFunctions}, car toute $f\in \sC$ est homéomorphe à une fonction $\hat{f}$ dans $\functionsonbaire$, et donc $f$ et $\hat{f}$ sont indistinguables par des propriétés topologiques qui sont celles qui nous intéressent. Cela éviterait certaines précisions dans la suite. Il faudrait aussi introduire $\functionsonbaire$ plus tôt, pour l'instant c'est dans le gluing. Once we establish that these scattered functions from a susbet of $\cN$ necessarily have countable (hence $0$-dim) range, we can simply say that $\sC$ is the class of scattered and continuous functions between zero-dimensional metrizable spaces. C'est plus joli, mais il nous faut la dérivée pour le voir.
%\end{aynote}
%
Here are few basic facts about our closed derivative.

\begin{fact}\label{CBbasics0}
Let $f$ be a function between topological spaces.
\begin{enumerate}
\item For every open set $U\subseteq\dom(f)$ and every ordinal $\alpha$ we have $\CB_\alpha(f\restr{U})=\CB_\alpha(f)\cap U$. \label{CBbasicsfromJSL2}
\item If $f$ scattered and $\CB(f)$ is limit, then for all $x\in\dom(f)$ there is an open $U\ni x$ with $\CB(f\restr{U})<\CB(f)$.
\item If $f$ is scattered and has nonempty compact domain, then $\CB(f)$ is a successor ordinal $\alpha+1$ and $f(\CB_\alpha(f))$ is finite.
\end{enumerate}
\end{fact}

\begin{proof}
\begin{enumerate}
\item {By induction on $\alpha$, observing that} for $U$ open and $x\in U$ we have $f\restr{U}$ is locally constant at $x$ if and only if $f$ is locally constant at $x$. 
\item As $\ker_{\CB}{f}=\bigcap_{\alpha<\CB(f)}\CB_\alpha(f)=\emptyset$, for all $x\in \dom f$ there exists $\alpha<\CB(f)$
such that $x$ belongs to the open set $U=\dom f \setminus \CB_\alpha(f)$. By the first item, it follows that $\CB(f\restr{U})\leq \alpha$. 
\item If $f$ had a limit $\CB$-rank $\lambda$, then the decreasing sequence of closed sets $(\CB_\alpha(f))_{\alpha<\lambda}$ would have empty intersection,
a contradiction with the compactness of $\dom f$.
So $\CB(f)=\alpha+1$ and $\CB_{\alpha}(f)$ is compact -- as a closed set of $\dom f$. Since $\CB_{\alpha+1}(f)=\emptyset$,
the sets $f^{-1}(\set{y})\cap \CB_{\alpha}(f)$ for $y\in \im(f\restr{\CB_\alpha(f)})$, form a partition of $\CB_\alpha(f)$ in nonempty open sets,
so $f$ must have must have finite range on $\CB_\alpha(f)$.\qedhere
\end{enumerate}
\end{proof}

The following basic results on the interplay between continuous reducibility and the $\CB$-derivative on functions is proven in \cite[Proposition 2.2 and Corollary 2.3]{carroy2013quasi} for Polish spaces,
but the same proofs work for arbitrary functions between topological spaces.

\begin{proposition}\label{CBbasicsfromJSL}
Let $f$ and $g$ be two functions between topological spaces. %Suppose that $g$ is scattered.
\begin{enumerate}
\item If $f\leq g$ and $g$ is scattered, then $f$ is scattered too. \label{CBbasicsfromJSL0}
\item If $(\sigma,\tau)$ continuously reduces $f$ to $g$, then for all $\alpha$ we have $\sigma(\CB_\alpha(f))\subseteq\CB_\alpha(g)$. \label{CBbasicsfromJSL1}
\end{enumerate}
\end{proposition}
\begin{proof}
Let $(\sigma,\tau)$ continuously reduce $f$ to $g$ and $A\subseteq \dom f$. We show that for all $x\in A$, if $g\restr{\sigma(A)}$ is locally constant at $\sigma(x)$, then $f\restr{A}$ is locally constant at $x$. If $g\restr{\sigma(A)}$ is locally constant at $\sigma(x)$, then there exists $U\ni \sigma(x)$ open in $\dom g$ such that $g(U\cap \sigma(A))$ is a singleton. By continuity of $\sigma$ there exists an open set $V\ni x$ of $\dom f$ such that $\sigma(V)\subseteq U$.
Hence $g\sigma(V\cap A)$ is a singleton, and therefore $f=\tau g \sigma$ is constant on $V\cap A$. This directly implies the first item and by induction it yields the second item.
\end{proof}

We now have all the elements to prove \cref{scatterediffemptykernel} from the introduction:
a continuous function between metrizable spaces is scattered if and only if it has non-empty perfect kernel if and only if $\id_\Q$ does not embed in it.

\begin{proof}[Proof of \cref{scatterediffemptykernel}]
After \cref{scatterediffemptykernel_general}, note that $\id_\Q$ is not scattered and use \cref{prop:nlc_implies_nonscattered} and \cref{CBbasicsfromJSL}~\cref{CBbasicsfromJSL0}.
\end{proof}

Here is useful relation on the $\CB$-rank of a scattered function and the ranks of its restrictions to open sets. %in the presence of an open cover of its domain.

\begin{corollary}\label{CBrankofclopenunion}
Let $f$ be a scattered function and $(A_i)_{i\in I}$ be an open covering of $\dom(f)$ for some set $I$.
Then $\CB(f)=\sup_{i\in I}\CB(f\restr{A_i})$.
\end{corollary}

\begin{proof}
Set $\alpha=\sup_{i\in I}\CB(f\restr{A_i})$. Using \cref{CBbasics0}~\cref{CBbasicsfromJSL2}, for all $i\in I$, we have $\CB_\alpha(f\restr{A_i})=\emptyset$ and 
 $\CB(f\restr{A_i})=\CB(f)\cap A_i$, and so: 
\[\CB_\alpha(f)\subseteq\bigcup_{i\in I}\CB_\alpha(f)\cap A_i=\bigcup_{i\in I}\CB_\alpha(f\restr{A_i})=\emptyset,\]
and so $\CB(f)\leq\alpha$. Now if $\beta<\alpha$, then $\beta<\CB(f\restr{A_i})$ for some ${i\in I}$ and
\[\CB_\beta(f)\supseteq\CB_\beta(f)\cap A_i=\CB_\beta(f\restr{A_i})\neq\emptyset.\]
Since $\CB_{\CB(f)}(f)=\emptyset$, it follows that $\beta<\CB(f)$, hence $\alpha=\CB(f)$.
\end{proof}

Given a scattered function $f$, we define as follows the \emph{$\CB$-degree}\index[IdxNotion]{Cantor--Bendixson!degree} of $f$ that we denote by $N_f$. When $\CB(f)$ is limit or null, set $N_f=0$.
When $\CB(f)=\alpha+1$ for some $\alpha<\omega_1$, we let $N_f$\index[IdxSymb]{N@{$N_f$}} be the cardinality of $f(\CB_\alpha(f))$. 
We also call \emph{$\CB$-type} of $f$ the pair $\tp(f)=(\CB(f),N_f)$. \index[IdxNotion]{Cantor--Bendixson!type}\index[IdxSymb]{tp@{$\tp(f)$}}

We denote by $\leq_{lex}$\index[IdxSymb]{2@{$\leq_{lex}$}} the lexicographical order on sequences of ordinals. Here is an immediate application of \cref{CBbasicsfromJSL} (see \cite[Fact, p. 642]{carroy2013quasi}).

\begin{corollary}\label{CBbasicsfromJSL-tp}
For any scattered functions $f$ and $g$, $f\leq g$ implies $\tp(f)\leq_{lex}\tp(g)$.
\end{corollary}
\begin{proof}
Assume that $f\leq g$ as witnessed by $(\sigma,\tau)$. By \cref{CBbasicsfromJSL}~\cref{CBbasicsfromJSL1}, we have $\sigma(\CB_{\CB(g)(f)})\subseteq \ker_{\CB} g=\emptyset$ so $\CB(f)\leq \CB(g)$. If $\CB(f)<\CB(g)$, then $\tp(f)\leq_{lex}\tp(g)$. So assume that $\CB(f)=\CB(g)=\alpha$. If $\alpha$ is limit, then by definition $N_f=N_g=0$. Otherwise, $\alpha=\beta+1$ and $\sigma(\CB_\beta(f))\subseteq \CB_\beta(g)$ by \cref{CBbasicsfromJSL}~\cref{CBbasicsfromJSL1}. 
Hence for every $x\in \CB_\beta(f)$ we have $f(x)=\tau ( g\sigma(x))$ with $g\sigma(x)\in g(\CB_\beta(g))$,
which shows that $\tau$ maps surjectively a subset of $g(\CB_\beta(g))$ onto $f(\CB_\beta(f))$, so $N_g \geq N_f$ as desired.
\end{proof}



\section{Disjoint union}


\begin{definition}
Given an index set $I$ and a sequence $(f_i)_{i\in I}$ of functions between topological spaces,\index[IdxNotion]{operation!disjoint union}
we say that a function $f$ is the \emph{disjoint union} of the maps $f_i$ and we write $f=\bigsqcup_{i\in I} f_i$\index[IdxSymb]{u@{$\bigsqcup_if_i$}}
if $f=\bigcup_{i\in I}f_i$ and $(\dom f_i)_{i\in I}$ is a clopen partition of $\dom(f)$.
\end{definition}

Note that the disjoint union over the empty set is the empty function.
The previous definition makes sense for any set $I$, but in this paper $I$ will always be countable, so in general we assume that $I\subseteq \N$.
When the context is clear, we write $(f_i)_i$ instead of $(f_i)_{i\in I}$.

The main reason for the disjoint union operation is its the following strong link with local properties for functions with $0$-dimensional domains.
Given a class $\mathcal{F}$ of functions, we say that a function $f$ is \emph{locally in $\mathcal{F}$}, or \emph{locally $\mathcal{F}$} for short,
if for all $x\in\dom(f)$ there is an open $U\ni x$ such that $f\restr{U}\in\mathcal{F}$.

Given sequences $s\in\N^{<\N}$ and $x\in\N^{\leq\N}$ we write $s\sqsubseteq x$\index[IdxSymb]{sq@{$s\sqsubseteq x$}} when $s$ is a \emph{prefix}, or an \emph{initial segment}, of $x$.
Recall that the sets $N_{s}=\set{x\in\cN}[s\segm x]$ for all $s\in\N^{<\N}$ form a basis of clopen sets for $\cN$.
A \emph{tree} $T$ (on $\N$) is a subset of $\N^{<\N}$ closed under prefixes, it is \emph{pruned} when all $s\in T$ have a strict extension $t\sqsupset s$ in $T$.
In this case we denote by $[T]$ the set of infinite branches of $T$.
Recall that sets of the form $[T]$ for pruned trees $T$ on $\N$ are in bijection with closed sets $C\subseteq\cN$ (see \cite[(2.4)]{kechris}).

\begin{proposition}\label{0dimanddisjointunion}
Let $f$ be a function with separable metrizable 0-dimensional domain and $\mathcal{F}$ a class of functions.
Then $f$ is locally $\mathcal{F}$ if and only if $f=\bigsqcup_{i\in I}f_i$ for some sequence of functions $(f_i)_{i\in I}\subseteq\mathcal{F}$.
\end{proposition}


\begin{proof}
Being separable metrizable and 0-dimensional, the domain of $f$ is homeomorphic to a subset of $\cN$,
so without loss of generality we can suppose that $X=\dom f$ is actually a subset of $\cN$.
Call $T$ the pruned subtree of $\N^{<\N}$ satisfying $[T]=\overline{X}$ the closure of $X$ in $\cN$. 
Let $S=\set{s\in T}[f\rest{(N_{s}\cap X)}\in\mathcal{F}]$ and define $B$ as the set of $\sqsubseteq$-minimal elements of $S$.
The elements of $B$ are pairwise incomparable for the prefix relation $\segm$, so the clopen sets $N_s$ for $s\in B$ are pairwise disjoint.
Since $f$ is locally $\mathcal{F}$, for all $x\in X$ there exists $s\in B$ with $s\sqsubset x$, so the sets $N_s\cap X$ for $s\in B$ form an open cover of $X$.
We have thus obtained a clopen partition $\set{N_s\cap X}[s\in B]$ of $\dom f$ such that $f\rest{(N_s\cap X)}\in\mathcal{F}$ for all $s\in B$, proving our point.
\end{proof}

Scattered functions $f$ with successor rank $\alpha+1$ which are constant on $\CB_\alpha(f)$ are ubiquitous and play an important role in the sequel. 
\begin{definition}
A function $f$ is said to be \emph{simple}\index[IdxNotion]{simple function}\index[IdxNotion]{function!simple} if it is scattered with $\CB$-degree $1$. A simple function has successor $\CB$-rank $\alpha+1$ for some ordinal $\alpha$ and we call the unique point in the image of $f$ on  $\CB_{\alpha}(f)$ the \emph{distinguished point of $f$}.
%We call \emph{simple} a function $f$ of $\CB$-degree $1$, note that the $\CB$-rank of a simple function is always successor.
%In this case, we call the unique point in the image of $f$ on  $\CB_{\CB(f)-1}(f)$ the \emph{distinguished point of $f$}.
\end{definition}
A first immediate but crucial application of \cref{0dimanddisjointunion} is the following essential lemma that generalizes \cite[Lemma 2.4]{carroy2013quasi} with the same proof.

\index[IdxNotion]{Decomposition Lemma}
\begin{lemma}[Decomposition Lemma]\label{JSLdecompositionlemma}
Any scattered function from a $0$-dim\-en\-sional separable metrizable space is locally simple.

If moreover it has successor rank $\beta+1$, then it is locally of type $(\beta+1,1)$. % it is moreover $\CB(f)=\beta+1$, then one can write $f=\bigsqcup_{i\in I}f_i$ with $\tp(f_i)=(\beta+1,1)$ for all $i\in I$.
\end{lemma}

\begin{proof}
Let $f:X\to Y$ be scattered with $X$ a $0$-dimensional separable metrizable.
By induction on $\CB(f)$. Vacuously true if $\CB(f)=0$. If $\CB(f)$ is limit, then every $x\in \dom f$ admits a clopen neighbourhood $C$ such that $\CB(f\restr{C})<\CB(f)$
and the result follows by induction hypothesis and \cref{0dimanddisjointunion}.
Finally, assume that $\CB(f)=\beta+1$ and let $I= f(\CB_\beta(f))$. Since $f$ is locally constant on the closed set $\CB_\beta(f)$,
we can choose for each $y\in I$ an open set $U_y$ of $\dom(f)$ such that $U_y\cap \CB_\beta(f)=f^{-1}(\{y\})\cap \CB_\beta(f)$.
Applying the generalized reduction property of open sets \cite[22.16]{kechris} to the open cover of $\dom(f)$ given by $V_y=U_y\cup (\dom(f)\setminus \CB_\beta(f))$ for $y\in I$
yields a clopen partition $(C_y)_{y\in I}$ of $\dom(f)$ with $C_y\subseteq V_y$ for all $y\in I$.
Note that for all $y\in I$ we have $C_y\cap \CB_\beta(f)= f^{-1}(\{y\})\cap \CB_\beta(f)$,
which readily implies that each $f\restr{C_y}$ is simple of $\CB$-rank equal to $\beta+1$ using \cref{CBbasics0}~\cref{CBbasicsfromJSL2}, as desired.
\end{proof}

\smallskip

We say that an operation on a class $\mathcal{C}$ of objects -- such as sets, spaces, functions -- \emph{preserves} a property $\mathcal{F}$ if\index[IdxNotion]{operation!preserving a property}
the operation gives a result in $\mathcal{F}$ whenever the operands are in $\mathcal{F}$.
For instance, the disjoint union operation preserves continuity by the following basic but useful criterion for continuity, similar to \cite[2.2.4 and 2.2.6]{engelking}.

\smallskip


Given a metrizable space $X$ and an index set $I$, we say that a family $(A_i)_{i\in I}$ of subsets of $X$ is a \emph{relative clopen partition}\index[IdxNotion]{relative clopen partition}\index[IdxNotion]{partition!relative clopen} in $X$
when $(A_i)_{i\in I}$ forms a clopen partition of the subspace $\bigcup_{i\in I} A_i$, or equivalently
if the sets $A_i$ are pairwise disjoint and relatively open in $\bigcup_iA_i$.
Note that a pairwise disjoint family $(A_i)_{i\in I}$ is a relative clopen partition exactly when the labelling function $c:\bigcup_i A_i \to I$,
$c(a)=i$ if and only if $a\in A_i$, is continuous when $I$ is endowed with the discrete topology.


\begin{lemma}\label{lem:ContUnion}
Let $X$ and $Y$ be topological spaces with $X$ metrizable and $(f_i:A_i\to B_i)_{i\in I}$ be countable sequence of continuous functions with $A_i\subseteq X$ and $B_i\subseteq Y$ for all $i\in I$.
If $(A_i)_i$ is a relative clopen partition, then $f=\bigsqcup_{i\in I} f_i$ is continuous.
\end{lemma}
\begin{proof}
Since $X$ is metrizable, it is enough to prove that for any sequence $(x_n)_{n\in\N}$ converging in $\bigcup_{i\in I}A_i$ to $x$, the sequence of images converges to $f(x)$.
If $(x_n)_{n\in\N}$ converges in $\bigcup_i A_i$, say to $x\in A_i$ for some $i$,
then as $A_i$ is open in $\bigcup_i A_i$, then for some $N\in\N$ we have $x_n\in A_i$ for all $n\geq N$. Since $(f_i(x_n))_{n\geq N}$ converges to $f_i(x)$ by continuity of $f_i$, it follows that $(f(x_n))_{n\in\N}$ converges to $f_i(x)=f(x)$.
\end{proof}

We remark that we can remove the metrizability assumption on $X$ in the above result by requiring that the sets $A_i$ are \emph{separated by open sets}, namely that
there is a family $(B_i)_ {i\in I}$ of pairwise disjoint open sets in $X$ such that $A_i\subseteq B_i$ for all $i\in I$. 
In fact, in metrizables spaces, relative clopen partitions coincide with families separated by open sets. This follows from the fact that metric spaces are paracompact~\cite{rudin1969new} and therefore hereditarily collectionwise normal~\cite[Theorem 5.1.18 and Problem 5.5.1]{engelking}.
For $0$-dimensional separable metrizable spaces, a direct proof of this fact can be given using similar ideas as in \cref{0dimanddisjointunion}.

%\rnote*{This has become marginal, we could just say that there is a direct proof in the zero-dimensional case with a similar idea as in \cref{0dimanddisjointunion}.}{
%In this paper, we focus on zero dimensional separable spaces and we now provide a direct proof of this equivalence for these class of spaces.
%
%\begin{fact}\label{FactRelCopen}
%Suppose that $X$ is a separable metrizable zero-dimensional space.
%A family $(A_i)_{i\in I}$ of subsets of $X$ is separated by open sets
%if and only if it is a relative clopen partition.
%\end{fact}
%\begin{proof}
%The very definition of separation by open sets gives the direct implication, it is enough to prove the converse.
%
%If $X$ is zero-dimensional metrizable separable, its completion is homeomorphic to a closed subset of $\cN$, so we can assume that the sets $A_i$ are subsets of $\cN$.
%Suppose that $(A_i)_{i\in I}$ is a pairwise disjoint family of sets, each of them relatively open in $\bigcup_{i\in I} A_i:=A$.
%For all $i$, let $F_i=\set{s\in\N^{<\N}}[N_s\cap A_i \neq \emptyset \text{ and } N_s\cap A_j=\emptyset \text{ for all $j\neq i$}]$. Let $B_i$ be the set of $\segm$-minimal elements of $F_i$ and set $U_i=\bigcup\set{N_{s}}[s\in F_i]$. Since $A_i$ is relatively open in $A$, for all $x\in A_i$ there exists $s\sqsubseteq x$ with $N_s$ disjoint from $\bigcup_{j\neq i}A_j$, so $A_i\subseteq U_i$.
% %We claim that the open sets $U_i=\bigcup\set[N_{s}]{s\in F_i}$ are as desired.
%%By definition, $A_i\subseteq U_i$ for all $i$. 
%Suppose that there exists $x\in U_i\cap U_j$ for some $i\neq j$, then there exists $s\in B_i$ and $t\in B_j$ that are comparable.
%Assume, without loss of generality, that $s\segm t$ holds.
%Since $t\in B_j$, there exists $y\in A_j$ such that $y\in N_t$ and so $y\in N_s$. However, $N_s\cap A_j=\emptyset$ because $s\in B_i$, a contradiction.
%\end{proof}
%}



\section{Gluing}\label{subsectionGluing}


We close this preliminary section with the study of the Gluing operation which provides a natural \enquote{addition} for functions. Combined with the Pointed Gluing operation, also defined in~\cite{carroy2013quasi} and studied in \cref{sectionPointedGluing}, this operation will allow us to define  maximum (respectively minimum) functions among functions in $\sC$ with $\CB$-rank at most $\alpha$ (respectively at least $\alpha$) for all countable ordinal $\alpha$. %, thanks to \cref{JSLdecompositionlemma}.
%whose hypotheses are almost - continuity excepted - those satisfied in $\sC$.
%\ynote*{pertinent?}{Even though our main focus is indeed on $\sC$, we try to state our results in more generality when it does not make them too cumbersome.}


Given sequences $s\in\N^{<\N}$ and $x\in\N^{\leq\N}$ we denote by $s\conc x$ the concatenation of $s$ with $x$,\index[IdxSymb]{conc@{$s\conc x$, $s\conc A$}}
and for any $A\subseteq\N^{\leq\N}$ the notation $s\conc A$ stands for $\set{s\conc x}[x\in A]$.%\index[IdxSymb]{$s\conc A$}
The following operation is called \emph{sum of the mappings} in \cite[2.2.E]{engelking}.

\begin{definition}
We call \emph{gluing}\index[IdxNotion]{operation!gluing} of a countable sequence $(A_i)_ {i\in I}$ of subsets of $\cN$ the set $\gl_{i\in I}A_i=\bigcup_{i\in I}(i)\conc A_i$.\index[IdxSymb]{op@{$\gl_iA_i$}}
Given a countable sequence of functions $(f_i:A_i\rao B_i)_{i\in I}$ in $\functionsonbaire$,
We call \emph{gluing} of $(f_i)_{i\in I}$, the map\index[IdxSymb]{op@{$\gl_if_i$}}
\begin{align*}
\gl_{i\in I}f_i : \gl_{i\in I}A_i &\longrightarrow\gl_{i\in I}B_i\\
 (i)\conc x &\longmapsto (i)\conc f_i(x).
\end{align*}
\end{definition}

The following facts are either classical exercises left to the reader, or a direct consequence of \cref{CBrankofclopenunion}.

\begin{Fact}\label{BasicsOnGluing}
\
\begin{enumerate}
\item The gluing operation preserves countability\footnote{Because we only consider \emph{countable} gluings here.}, openness, closure, and Polishness of sets; continuity, injectivity, surjectivity and scatteredness of functions.
\item For $I\subseteq \N$ and $(f_i)_{i\in I}$ in $\functionsonbaire$, we have $\CB(\gl_if_i)=\CB(\bigsqcup_if_i)=\sup_i\CB(f_i)$.\index[IdxNotion]{Cantor--Bendixson!rank!of a gluing}
\item Gluing \emph{commutes with the identity natural transformation}, that is for any $(X_i)_{i\in I}\subseteq \cN$ we have $\id_{\gl_{i\in I}X_i}=\gl_{i\in I}(\id_{X_i})$.
\end{enumerate}
\end{Fact}

\subsection{Gluing as an upper bound}

Our main task for each of our three operations is to prove sufficient -- and, whenever possible, necessary -- conditions on operands $(g_i)_i$
for the result of the operation to reduce (or be reduced by) an arbitrary function $f$.
The next result represents the first step, characterizing the situation where $\gl_i g_i$ is an upper bound for a given function $f$.
These results are all new, except for the next one that was somehow already present in \cite[Lemma 3.3]{carroy2013quasi}.
%\rnote{j'ai tout commenté car je trouvais ça pas clair...}
%\ynote*{modifié}{via restrictions of $f$ on a specific partition}. %the existence of a specific partition of its domain.

\index[IdxNotion]{upper bound condition!gluing}\index[IdxNotion]{gluing operation!upper bound condition}
\begin{proposition}[Gluing as upper bound]\label{Gluingasupperbound}
Let $(g_i:X_i\to Y_i)_{i\in I}$ be countable sequence in $\functionsonbaire$. %, $I\subseteq\N$, and for all $i\in I$ $g_i:X_i\to Y_i$ be
The following are equivalent for every function $f:A\to B$ between topological spaces:
\begin{enumerate} 
\item $f\leq \gl_{i\in I}g_i$, %there is a continuous reduction from $f$ to $\gl_{i\in I}g_i$ if and only if
\item there is a clopen partition $(A_i)_{i\in I}$ of $A$ such that $f\restr{A_i}\leq g_i$ for all $i\in I$.
\end{enumerate}
\end{proposition}

\begin{proof}
Suppose that for all $i\in I$ there is a continuous reduction $(\sigma_i,\tau_i)$ from $f\restr{A_i}$ to $g_i$. Set then $\sigma:x\mapsto (i)\conc \sigma_i(x)$ if $x\in A_i$ and $\tau:(i)\conc y\mapsto\tau_i(y)$. The map $\tau$ is continuous by \cref{lem:ContUnion}. Since the sets $A_i$ are pairwise disjoint, $\sigma$ is well-defined, and it is continuous because the sets $A_i$ are clopen.
The fact that $(\sigma_i,\tau_i)$ is a reduction for all $i\in I$ implies that $(\sigma, \tau)$ is a reduction from $f$ to $\gl_{i\in I}g_i$.

Suppose now that $(\sigma,\tau)$ continuously reduces $f$ to $\gl_i g_i$, and set $A_i=\sigma^{-1}(N_{(i)})$.
Since $(N_{(i)})_i$ is a clopen partition of $\dom(\gl_{i\in I} g_i)$ and $\sigma$ is continuous, $(A_i)_i $ is a clopen partition of $\dom(f)$.
Conclude by observing that $f\restr{A_j}\leq(\gl_{i\in I}g_i)\restr{N_{(j)}}\equiv g_j$ for all $j\in I$.
\end{proof}

The following simple corollary will be used extensively in the sequel.

\begin{corollary}\label{Gluingasupperbound_cor}
Let $f:A\to B$ in $\functionsonbaire$ and let $\partit$ be a clopen partition of $A$.
Then $f=\bigsqcup_{P\in\partit} f\restr{P}\leq \gl_{P\in\partit} f\restr{P}$.
%For $I\subseteq\N$, let $(f_i:A_i\to B_i)_{i\in I}$ be in $\functionsonbaire$.
%If $(A_i)_{i\in I}$ is a relative clopen partition, then $\bigsqcup_{i\in I}f_i\leq\gl_{i\in I} f_i$.
\end{corollary}

\subsection{Gluing as a lower bound}
Dually, bounding a function $f$ from below with a gluing involves corestrictions of $f$.

\index[IdxNotion]{lower bound condition!gluing}\index[IdxNotion]{gluing operation!lower bound condition}
\begin{Proposition}[Gluing as lower bound]\label{Gluingaslowerbound}
Let $(g_i:X_i\to Y_i)_{i\in I}$ be a countable sequence in $\functionsonbaire$.
%$f:A\to B$ be a function, $I\subseteq\N$,
%and $g_i:X_i\to Y_i$ be in $\functionsonbaire$ for all $i\in I$.
%Then for every function $f:A\to B$ between metrizable spaces, $\gl_{i\in I}g_i\leq f$ if and only if there is a family $(B_i)_{i\in I}$ of subsets of $B$ satisfying:
%\begin{enumerate}
%\item the family $(B_i)_{i\in I}$ is a relative clopen partition,
%\item for all $i\in I$ we have $g_i\leq f\corestr{B_i}$.
%\end{enumerate}
The following are equivalent for every function $f:A\to B$ between metrizable spaces:
\begin{enumerate}
\item $\gl_{i\in I}g_i\leq f$
%\Needspace{4\baselineskip}
\item there exists a family $(U_i)_{i\in I}$ of subsets of $B$ satisfying:
\begin{enumerate}
\item $(U_i)_{i\in I}$ is a relative clopen partition, and
\item for all $i\in I$ we have $g_i\leq f\corestr{U_i}$.
\end{enumerate}
\end{enumerate}
\end{Proposition}

\begin{proof}
Suppose first that $(\sigma,\tau)$ continuously reduces $\gl_{i\in I}g_i$ to $f$, then the sets $U_i=\tau^{-1}(N_{(i)})$, $i\in I$, 
form a clopen partition of $\dom \tau=\im f\sigma$. 
As $g_i\leq(\gl_{j\in I}g_j)\corestr{N_{(i)}}=(\tau f\sigma)\corestr{N_{(i)}}=(f\sigma)\corestr{U_i}\leq f\corestr{U_i}$, we have the direct implication.

For all $i\in I$, we fix a continuous reduction $(\sigma_i,\tau_i)$ from $g_i$ to $f\corestr{U_i}$,
and set $\sigma:(i)\conc x\mapsto\sigma_i(x)$ for $x\in X_i$ and $\tau: y\mapsto (i)\conc \tau_i(y)$ for $y\in \dom \tau_i$.
Since the sets $U_i$ are pairwise disjoint and $\dom \tau_i\subseteq U_i$, $(\sigma,\tau)$ is a well-defined pair of maps. Note that $\sigma$ is continuous by \cref{lem:ContUnion}.
Since $(U_i)_i$ is a relative clopen partition and $\dom \tau_i\subseteq U_i$, so is the family $(\dom \tau_i)_i$ and $\tau$ is continuous by \cref{lem:ContUnion} too.
The fact that $(\sigma_i,\tau_i)$ are reductions for all $i\in I$ implies that $(\sigma, \tau)$ is a reduction from $\gl_{i\in I}g_i$ to $f$.
Indeed, for $x\in X_i$ we have $\tau f \sigma ((i)\conc x)=(i)\conc \tau_i f \sigma_i(x)= (i)\conc g_i(x)$.
\end{proof}

We use the following simple corollary extensively in the sequel.

\begin{corollary}\label{Gluingaslowerbound2}
Let $f:A\to B$ be in $\functionsonbaire$. If $\mathcal{U}$ is a family of pairwise disjoint open sets of $B$, then $\gl_{U\in\mathcal{U}}f\corestr{U}\leq f$.
%If $(B_i)_{i\in I}$ is a relative clopen partition in $\im f$ then $\gl_{i\in I}f\corestr{B_i}\leq f$.
\end{corollary}

In order to prove that a function $f$ is continuously equivalent to a gluing of functions, one needs to consider both restrictions and corestrictions of $f$. 
One needs to find both a clopen partition of the domain of $f$ to use \cref{Gluingasupperbound} and a relatively clopen partition in the image of $f$
to use \cref{Gluingaslowerbound}.
While a clopen partition of the image of a continuous $f$ induces a clopen partition of the domain as in the following corollary,
we will also consider families in the domain and the range of $f$ that are not related in this way (see \cref{SolvableDecomposition}).

\begin{corollary}\label{UsefulcriterionforequivFinGl}
Let $f:A\to B$ be a continuous function between topological spaces
and $(g_i:X_i\to Y_i)_{i\in I}$ be countable sequence in $\functionsonbaire$.

If there exists a clopen partition $B=\bigsqcup_{i\in I} B_i$ with $f\corestr{B_i}\equiv g_i$ for all $i\in I$, then $f\equiv \gl_{i\in I} g_i$.
\end{corollary}

For any countable ordinal $\alpha$, we write $\sC_\alpha$ (resp. $\sC_{\leq \alpha}, \sC_{<\alpha}$) for the set of all functions of $\CB$-rank $\alpha$ (resp. $\leq\alpha,<\alpha$) in $\sC$.
\index[IdxSymb]{Scat@{$\sC_\alpha$, $\sC_{\leq\alpha}$, $\sC_{<\alpha}$}}

For instance, $\sC_0$ consists only of the empty function, and $\sC_1$ is the class of locally constant functions.

A function of finite $\CB$-degree has finite image on the last $\CB$-derivative, so the previous corollary immediately gives the generalization to $\sC$ of \cite[Lemma 4.3]{carroy2013quasi}:

\begin{corollary}\label{FiniteDegreeAreFinGl}
Let $\alpha<\omega_1$ and $f\in\sC_{\alpha+1}$.
If the $\CB$-degree of $f$ is $n+1$, then $f\equiv\gl_{i\leq n}f_i$ for some sequence $(f_i)_{i\leq n}$ of simple functions in $\sC_{\alpha+1}$.
\end{corollary}
\begin{proof}
Take a finite clopen partition of the co-domain of $f$ such that each piece contains exactly one point of the image of $f$ on $\CB_\alpha(f)$ and apply the previous corollary.
\end{proof}


Given functions $f_0$ and $f_1$ in $\functionsonbaire$, we write $f_0\glbin f_1$ instead of $\gl_{i\in 2}f_i$ for the binary gluing operation.
For $n\leq \omega$ and $f\in \functionsonbaire$, we write $nf$ for $\gl_{i\in n} f$.\index[IdxSymb]{op@{$f_0\glbin f_1$}}\index[IdxSymb]{op@{$nf$}}
A direct application of \cref{Gluingasupperbound} gives:

\begin{fact}\label{Gluingcohomomorphism}~
Let $(f_i)_{i\in I}, (g_k)_{k\in K} \subseteq \functionsonbaire$. If there exists $\iota:I\to K$ with $f_i\leq g_{\iota(i)}$ for all $i\in I$, then $\gl_{i\in I}f_i\leq \gl_{k\in K} g_k$. 

In particular, if $f\in\functionsonbaire$ and $m\leq n\leq \omega$, then $mf\leq nf$.
\end{fact}

\begin{definition}
If $\mathcal{F}$ is a set of functions, then we denote by $\FinGl{\mathcal{F}}$\index[IdxSymb]{Fin@{$\FinGl{\mathcal{F}}$}} the set of all finite gluings $\gl_{i< n} f_i$ for some finite sequence $(f_i)_{i<n}$ of functions in $\mathcal{F}$.
We say that a set $\mathcal{C}$ of functions is \emph{generated} by a set $\mathcal{F}$ if for all $g\in\mathcal{C}$ there exists $f\in \FinGl{\mathcal{F}}$ such that $g\equiv f$.
\index[IdxNotion]{finitely generated set of functions}

We say that $\mathcal{C}$ is \emph{finitely generated} if it generated by a finite set. 
\end{definition}

Remark that for $\mathcal{C}$ to be generated by $\mathcal{F}$, we do not require $\mathcal{F}$ to be a subset of $\mathcal{C}$. 
Note the empty function always belongs to $\FinGl{\mathcal{F}}$ as the gluing of the empty sequence and every element of $\FinGl{\mathcal{F}}$ takes the form $\gl_{i\leq k} n_i f_i$ where $k\in\N$, $n_0,\ldots,n_k\in \N$ and $f_0,\ldots,f_k\in \mathcal{F}$ are pairwise distinct. 

Given two quasi-orders $(P,\leq_P)$ and $(Q,\leq_Q)$, we say that a map $\varphi:P\to Q$ is a \emph{co-homomorphism}\index[IdxNotion]{co-homomorphism} when for all $p,p'$ in $P$ we have
$\varphi(p)\leq_Q\varphi(p')$ implies $p\leq_P p'$. In such a case, if $Q$ is a \wqo{} (resp. a \bqo{}) then so is $P$ (\cite[Lemma 3.10]{wellbetterinbetween}).

Recall also that \wqo{}s (resp. \bqo{}s) are stable by finite products, and that all well-orders are \bqo{}s (\cite[Example 3.9 and Proposition 3.14]{wellbetterinbetween}).
\index[IdxNotion]{better-quasi-order (\bqo{})}

\begin{proposition}\label{SecondstepforBQOthm}
Continuous reducibility is a \bqo{} on any finitely generated set of functions.
\end{proposition}

\begin{proof}%\ynote{it may seem from this proof that we need AC to define the co-homorphism, but in fact the proof that $\mathcal{C}$ is \bqo{} is by contradiction and only requires countable choice. If there was bad multi-sequence $g$ in $\mathcal{C}$, we can choose for ever function $f$ in the (countable) range of $g$ an $s_f\in \N^{k+1}$, then composing $g$ with $f\mapsto s_f$ yields a bad multi-sequence in $\N^{k+1}$.}
Let $\mathcal{C}$ be generated by a finite set of functions $\mathcal{F}=\set{f_0,\ldots, f_k}$ for some $k\in\N$.
For all $f\in\mathcal{C}$ there is an $s_f\in\N^{k+1}$ such that $f\equiv \gl_{i\leq k}s_f(i)f_i$.
But then if $f,g$ in $\mathcal{C}$ satisfy $s_f(i)\leq s_g(i)$ for all $i\leq k$, then $f\leq g$ by \cref{Gluingcohomomorphism},
so $f\mapsto s_f$ is a co-homomorphism from $\mathcal{C}$
to the finite product of well-orders $\N^{k+1}$, which makes $\mathcal{C}$ a \bqo{}.
\end{proof}


The previous proposition motivates our strategy underlined in \cref{Introthm:LevelsAreFinitelyGenerated}.
To prove that each level $\sC_\alpha$ is a \bqo{}, we actually describe a finite set of functions that generates it in the above sense.
The level $\sC_0$ consists only of the empty function, so
the first interesting level is $\sC_1$, that of locally constant functions, to which we now turn.

We use several time in this paper the following well-known exercise.

\begin{fact}\label{InfiniteEmbedOmega}
Any infinite metrizable space contains an infinite discrete subspace.
\end{fact}

We close this section with the analysis of the locally constant functions up to continuous equivalence.

\index[IdxNotion]{locally constant function}\index[IdxNotion]{function!locally constant}
\begin{proposition}\label{LocallyConstantFunctions}
The class of locally constant functions $f:X\to Y$ from a second countable space $X$ to a metrizable space $Y$ is (finitely) generated by\footnote{For notational simplicity we adopt the convention that $1=\set{\iw{0}}$.} $\set{\id_{1},\id_\N}$. In fact, $f\equiv \id_{I}$ where $I$ is a discrete set with cardinality $|\im f|$. 
\end{proposition}
\begin{proof}
Observe first that any constant function on a non-empty topological space is continuously equivalent to $\id_{1}$.
Now if $f$ is locally constant, the sets $A_y=f^{-1}(\set{y})$ for $y\in\im(f)$ form a clopen partition of $\dom(f)$ such that $f\restr{A_y}\equiv \id_{1}$ for all $y\in \im f$.
By \cref{Gluingasupperbound}, we get $f\leq \gl_{y\in \im f} \id_{1}$ and note that since $X$ is second countable, $\im f$ is countable. 

If $\im f$ is finite of cardinality $n$, then it is discrete and so using \cref{Gluingaslowerbound} with $B_y=\{y\}$ for $y\in \im f$,
we get $n \id_{1} \leq f$. Therefore $f\equiv n  \id_{1}$ in this case.

If $\im f$ is infinite, being metrizable it contains a countably infinite discrete subspace by \cref{InfiniteEmbedOmega},
so we similarly obtain $\gl_{n\in\N} \id_{1} \leq f$ using \cref{Gluingaslowerbound}.
Therefore in this case $f\equiv \gl_{n\in\N} \id_{1} \equiv \id_{\N}$, which concludes the proof.
\end{proof}

%\begin{aynote}{Remark. Ça m'a frappé à quel point le gros de notre travail consiste à généraliser la preuve du cas localement constant. Il me semble que souligner cela peut aider le lecteur à comprendre ce qu'on veut faire. Je me suis essayé. Peut-être que ça peut faire une petite pause bienvenue dans l'exposition.}
In the sequel we introduce the necessary concepts which allows us to generalize the previous simple proof to all successor levels of $\sC$.
Crucially, we define \emph{centered functions} which generalizes constant functions in \cref{sectionCentered}.
Similarly to the unique image of a constant function, centered functions admit a special point in their image, their \emph{\cocenter{}}.
While there is, up to equivalence, only one constant function, we will show that there are, up to equivalence, only finitely many centered functions of a given $\CB$-rank 
\cref{finitenessofcenteredfunctions}.
We eventually prove in \cref{PreciseStructureCor1} that every function $f$ is locally centered.
%, when continuous reducibility is \bqo{} on functions strictly below $f$. While this gives us a partition of the domain by \cref{0dimanddisjointunion}, note that in the above proof we gathered all constant parts of $f$ with same image $y$ in $f\restr{A_y}$. Here $f\restr{A_y}$ is still equivalent to the constant function $\id_{\set{0}}$, however gathering several equivalent centered parts with same \emph{cocenter} leads to a more complex situation, which is dealt with in \cref{VerticalTheorem}. Just like we got $f\leq  \gl_{y\in \im f} \id_{\set{0}}=G$ above, with a more work we will be able to upper bound $f$ with the right gluing $G$ of functions constructed from centered parts of $f$. Note that to obtain the other inequality $G\leq f$ in the case of locally constant functions, we pick an infinite discrete set in the range of $f$ which allows us to use \cref{Gluingaslowerbound}. %Note that this is required to upper bound the gluing $G$ that we identified in the first step and obtain the desired equivalence. 
%A somewhat similar -- though more elaborate -- selection of cocenters will be performed notably in \cref{SolvableDecomposition}. 

%This sketch inspired from the case of locally constant functions aims at proving finite generation for successor levels of $\sC$.
In the next section, we study the general structure of $\sC$ which mainly establishes that the structure of limit levels is as simple as one could hope for and ensures that finite generation of all levels indeed entail that $\sC$ is \bqo{}.
%\end{aynote}

%\begin{arnote}{à ce sujet}
%Je suis tout à fait en faveur d'une petite pause dans l'exposition à ce stade!
%Cela étant, je trouve ton premier paragraphe trop difficile à lire pour quelqu'un qui n'a pas encore lu la section 4.
%En plus, je ne suis moi-même pas totalement certain de comprendre l'intuition forte que tu as sur la généralisation du cas constant...
%On en a déjà parlé, tu as remarqué quelque chose ici qui je dois l'avouer m'échappe encore!
%
%Je vais essayer de penser à une façon de reprendre ce que tu as fait, peut-être en plus synthétique. J'y reviendrai.
%\end{arnote}
